% MATH companion editorial bridge VI/26; explanatory prose only.
% Avoided: new problem environments, labels, changed mathematical source.
\paragraph{How finite-type hypotheses enter the geometry.}
The phrase \emph{finite type} measures how many algebraic generators are needed
\emph{locally over the base}.  In an affine chart a morphism
\(\operatorname{Spec}B\to\operatorname{Spec}A\) is of finite type precisely when
\(B\) is generated as an \(A\)-algebra by finitely many elements.
In geometric terms, a finite list of coordinate functions describes points
in the fibers, although the equations satisfied by those coordinates may
still require separate attention if the base ring is not Noetherian.

For an arbitrary morphism \(f:X\to Y\), the adjective \emph{locally of finite
type} means that the affine algebra condition holds after taking suitable
open affine charts in both schemes.  To pass from locally of finite type to
finite type, one also requires \emph{quasi-compactness} of \(f\).
This is a condition on the entire inverse image of quasi-compact opens,
not merely on the number of generators within a single chart.

\paragraph{First comparison: finite type is not finite.}
Let \(k\) be a field and set
\[
 A=k[t],\qquad B=A[u]/(tu-1)\simeq k[t,t^{-1}].
\]
The associated morphism
\(\operatorname{Spec}B\to\operatorname{Spec}A\)
is an open immersion onto \(D(t)\).
It is of finite type because \(B\) has one algebra generator \(u\),
subject to \(tu=1\).  Nevertheless, \(B\) is not a finite \(A\)-module:
no finite collection of powers of \(t^{-1}\) generates all negative powers
under multiplication by polynomials in \(t\).
Thus \emph{finite type} must not be confused with a \emph{finite morphism}.

\paragraph{Second comparison: locally of finite type need not be finite type.}
Consider the morphism
\[
 \coprod_{n\geq 1}\operatorname{Spec}k
 \longrightarrow \operatorname{Spec}k.
\]
Every component is an affine scheme of finite type over \(k\), so the
morphism is locally of finite type.  Its source is not quasi-compact:
the components form an open cover with no finite subcover.
Consequently the morphism is not of finite type.  The example isolates
exactly the global finiteness condition that local algebra alone cannot
supply.

\paragraph{Why Noetherianity matters in the standard proof.}
Noetherianity is a restriction on ideals: every ideal is finitely
generated.  For a Noetherian base ring \(A\), the Hilbert basis theorem
implies that a finitely generated \(A\)-algebra \(B\) is Noetherian.
This is the reason finite-type morphisms over locally Noetherian schemes
have locally Noetherian sources.
Without the Noetherian hypothesis, finite generation as an algebra
does not automatically yield finite generation of every ideal.
The local-to-global arguments in the following problems use a finite
standard-open cover to transfer these ring-theoretic properties without
silently assuming that a non-quasi-compact cover is finite.
